\chapter{Introduction}\label{cha:introduction}

\section{Background}

\subsection{Scale and Criticality of Contemporary IoT Deployment}

The Internet of Things has moved from an emerging technology to load-bearing digital infrastructure. IoT Analytics reports that the global installed base of connected IoT devices reached 18.5 billion at the end of 2024, a 12\% increase over 2023, and projects growth of a further 14\% to approximately 21.1 billion by the end of 2025, rising to an estimated 39 billion by 2030 at a compound annual growth rate of 13.2\% \cite{IoTAnalytics2025}. These figures deliberately exclude general-purpose computing devices such as laptops, smartphones and tablets, and therefore represent dedicated ``things'': industrial sensors, network video recorders, point-of-sale terminals, medical monitors, gateways and controllers \cite{IoTAnalytics2025}.

The significance of this expansion is not merely one of volume. IoT endpoints are increasingly embedded in systems where availability and integrity carry physical, rather than purely informational, consequences: industrial control loops, municipal utility telemetry, and clinical monitoring. A compromise in these settings is not confined to data loss; it can interrupt a process, degrade a safety function, or disable a service on which people depend. Consequently, the security posture of the endpoint population has become a question of operational resilience rather than of IT hygiene alone.

This shift is visible across sectors that have historically been treated as separate concerns by security research but that now share a common endpoint profile. Hospitals depend on networked infusion pumps, patient monitors and building-management sensors that were, in the great majority of cases, never designed to defend themselves against a network-borne adversary; water and power utilities depend on remote telemetry and control points scattered across a geographically dispersed footprint that cannot be physically hardened at scale; and manufacturing and logistics operators depend on programmable controllers and sensor meshes that must run continuously, often for a decade or longer, without a maintenance window in which security software could be patched or replaced. What unites these otherwise dissimilar sectors is that the endpoint, not the perimeter, is where risk now concentrates, and the endpoint is precisely the tier this project addresses.

\subsection{Resource Constraints and the Endpoint Security Gap}

The defining technical characteristic of IoT endpoints, and the root of their security problem, is severe resource constraint. RFC 7228 formalises this by defining three classes of constrained devices: Class 0 (much less than 10 KiB of RAM and much less than 100 KiB of code storage), Class 1 (approximately 10 KiB RAM and 100 KiB code storage), and Class 2 (approximately 50 KiB RAM and 250 KiB code storage) \cite{Bormann2014}. The same document observes that these boundaries move only slowly, because in the embedded space the gains from increased transistor density are typically reinvested in lower unit cost and lower power draw rather than in additional computing capability \cite{Bormann2014}.

This has a direct architectural consequence. Conventional endpoint security agents-host-based intrusion detection, endpoint detection and response, signature databases, and full cryptographic stacks-assume memory and processing budgets that are orders of magnitude beyond what these classes provide. Security must therefore be relocated: either to the network edge, where a gateway with modestly greater resources can observe device traffic passively, or to the cloud, where computation is effectively unbounded but latency and bandwidth are not. Neither location is sufficient in isolation, and the tension between them motivates the two-tier design presented in this report.

The constraint is not purely a matter of memory footprint; it is also a matter of energy and duty cycle. Many of the devices under discussion are battery-powered or energy-harvesting, and every millijoule spent on a cryptographic handshake, a signature update, or a background monitoring agent is a millijoule not spent on the device's primary sensing or actuation function. This creates a persistent incentive, at the manufacturing stage, to minimise on-device security logic in favour of battery life and unit cost, which in turn means that whatever security capability a fleet of such devices ends up with is largely determined by what can be added \emph{around} the device rather than \emph{on} it. That observation is the architectural premise underlying the edge-and-cloud split adopted in this project: rather than asking a Class 1 or Class 2 endpoint to run an anomaly detector, the detector is pushed one hop away, to a gateway that can absorb the computational and energy cost on the device's behalf.

The consequences of this gap are measurable. Forescout's fifth annual analysis of connected-device risk reported a 15\% year-over-year rise in average device risk score, from 7.73 in 2024 to 8.98 in 2025, with routers alone accounting for more than half of all devices carrying the most critical vulnerabilities \cite{Forescout2025}. The device types identified as riskiest within the IoT category-network video recorders, IP cameras, VoIP equipment, network-attached storage and point-of-sale systems-are, notably, long-standing rather than novel problems \cite{Forescout2025}. This persistence is itself informative: it suggests that the gap is not a temporary consequence of immature tooling that the market will eventually resolve, but a structural feature of how constrained devices are designed, procured and deployed, which is unlikely to close without a change in where the defensive workload is placed.

\subsection{The Evolving IoT Threat Landscape}

The canonical demonstration of aggregated IoT weakness remains the Mirai botnet. Antonakakis et al. document its growth to a peak of roughly 600,000 concurrent infections and its use against Krebs on Security, OVH and Dyn during late 2016; the initial attack on Krebs on Security exceeded 600 Gbps, sourced from what the authors describe as some of the least powerful hosts on the Internet \cite{Antonakakis2017}.

That capability has not diminished. In its Q3 2025 DDoS threat report, Cloudflare attributed a record attack peaking at 29.7 Tbps and 14.1 billion packets per second, lasting 69 seconds, to the Aisuru botnet-a TurboMirai-class IoT botnet estimated to comprise between one and four million compromised hosts, principally routers, CCTV cameras and DVR systems \cite{Cloudflare2025Q3, Kovacs2025}. Cloudflare reported mitigating 2,867 Aisuru attacks during 2025, of which 1,304 hyper-volumetric events occurred in Q3 alone \cite{Cloudflare2025Q3}. A subsequent attack peaking at 31.4 Tbps and lasting 35 seconds was reported in its Q4 2025 analysis \cite{Cloudflare2026Q4}. Comparing the 600 Gbps of 2016 with the tens of terabits per second of 2025 indicates roughly two orders of magnitude of growth in attainable attack volume drawn from the same class of endpoint.

Two properties of this traffic are directly relevant to defensive design. First, the attacks are short: Cloudflare observed that 71\% of HTTP-layer and 89\% of network-layer DDoS attacks concluded in under ten minutes \cite{Cloudflare2025Q3}. An incident that begins and ends inside ten minutes cannot be contained by a workflow that requires an analyst to receive an alert, interpret it, and manually configure a mitigation. Second, the compromised population is composed of devices that are, individually, unremarkable-which is precisely why per-device behavioural monitoring, rather than perimeter signature matching, is required to identify recruitment early.

The threat is not confined to volumetric botnet activity aimed outward at third parties; the same endpoint population is also the direct target of intrusions aimed at the operators who own it. An industry survey of 350 water and electricity operators across the United States and United Kingdom found that 62\% had been targeted by a cyberattack in the preceding year, and that 80\% of those had been targeted more than once; 54\% reported permanent corruption or destruction of data and operational systems as a result \cite{Semperis2025}. The same study attributed close to 60\% of the observed attacks to nation-state actors, and reported that compromised identity infrastructure, rather than the field devices themselves, was frequently the entry point-underscoring that the IoT and operational-technology layers of these organisations do not fail in isolation from the rest of the network, but are both a target and a launch point within a single, interconnected attack surface \cite{Semperis2025}. For a defensive architecture aimed at this endpoint population, this means that success cannot be judged solely by outward-facing metrics such as botnet participation; sustained, repeated targeting of the same device population, sometimes by well-resourced adversaries, must be treated as the baseline threat condition rather than the exceptional case.

\subsection{Limitations of Prevailing Detection and Response Practice}

Even where detection capability exists, the time to act on it has not improved. Mandiant's M-Trends 2025 reported a global median dwell time-the interval between initial compromise and detection-of 11 days for 2024, up from 10 days in 2023, against 205 days a decade earlier in 2014 \cite{Mandiant2025}. The following year's report recorded a further increase to 14 days for 2025, attributed in part to adversaries favouring edge devices such as routers, VPN appliances and firewalls, which typically ship without the logging depth that security teams expect from servers and endpoints \cite{Mandiant2026}. In other words, the very device class this project addresses is identified as a contributing blind spot.

The financial and operational consequences of this lag are quantifiable and, in physically consequential sectors, disproportionately severe. IBM's 2025 Cost of a Data Breach Report placed the global average cost of a breach at USD 4.44 million, the first year-on-year decline in five years, driven largely by faster containment through automated defences \cite{IBM2025}. Yet the same report found that healthcare organisations, which rely heavily on networked monitors, infusion pumps and building-management IoT, recorded the highest breach cost of any sector for the fifteenth consecutive year, with a breach lifecycle of 279 days-substantially longer than the 241-day global average-meaning that the sector with the greatest physical stakes is also the slowest to detect and contain an incident \cite{IBM2025}. This pattern is consistent with the broader argument of this report: where the consequences of compromise are most severe, the prevailing detection-to-response pipeline is least equipped to act quickly, which is the gap that motivates pushing detection and an initial remediation response down to the edge rather than depending exclusively on a centralised, human-mediated investigation.

Compounding this is the well-documented problem of alert fatigue, itself aggravated by a persistent shortage of qualified personnel to triage the alerts that detection systems do produce. Tariq et al., in a systematic review of the literature and industry practice, characterise alert fatigue as a socio-technical failure arising from alert volume, false-positive burden, tool fragmentation and analyst workload, and note survey evidence that a substantial majority of practitioners report burnout and increasing workload \cite{Tariq2025}. The 2025 SANS Detection and Response Survey correspondingly found that 73\% of organisations rank false positives as their foremost detection challenge \cite{SANS2025}. This is not simply a matter of insufficient headcount: the 2025 ISC2 Cybersecurity Workforce Study, drawing on responses from over sixteen thousand practitioners, found that only 34\% of organisations reported having the right level of cybersecurity staffing, and that budget constraints-rather than a lack of available candidates-had for the first time overtaken skills shortages as the primary driver of understaffed security teams \cite{ISC22025}. A detection system that raises an alarm on every transient traffic burst therefore does not merely produce noise; it actively erodes the trust and attention capacity of an already-stretched analyst population on which subsequent genuine alerts depend. This observation directly motivates the rolling, decaying risk score developed in Chapter \ref{cha:methodology}, and, more broadly, the decision to automate not only detection but an initial round of remediation, so that the framework's effectiveness does not depend on the availability of a human analyst at the moment an attack begins.

\subsection{Regulatory Pressure Toward Secure and Resilient Connected Products}

The policy environment has begun to codify what industry practice has not achieved voluntarily. The United Kingdom's consumer connectable product security regime, comprising Part 1 of the Product Security and Telecommunications Infrastructure Act 2022 together with the associated 2023 Regulations, came into force on 29 April 2024 and prohibits universal or easily guessable default passwords, requires a published point of contact for vulnerability reporting, and requires manufacturers to declare a minimum security update period \cite{UKPSTI2023}. At EU level, Regulation (EU) 2024/2847-the Cyber Resilience Act-entered into force on 10 December 2024, with vulnerability and incident reporting obligations applying from 11 September 2026 and full applicability from 11 December 2027; penalties for breaches of the essential requirements reach EUR 15 million or 2.5\% of worldwide annual turnover \cite{EU2024CRA}. In the United States, NIST IR 8259 and its companion core baseline set out foundational cybersecurity activities and device capabilities expected of IoT manufacturers \cite{Fagan2020}, complemented internationally by the ETSI EN 303 645 baseline for consumer IoT \cite{ETSI303645}.

Taken together, these instruments share a common structure: they legislate what a device must be capable of at the point of manufacture-strong authentication, a disclosed vulnerability-reporting channel, a defined update period-rather than how a device, or the network it sits within, should behave once that device has already been compromised. This is a deliberate and reasonable scope for product-safety regulation, but it leaves an operational gap precisely where this project is positioned. A newly manufactured device that fully satisfies the Cyber Resilience Act's essential requirements can still be recruited into a botnet once deployed, and none of the reviewed instruments mandate an automated detection-and-remediation capability at the network or gateway level. In this sense, regulation and the framework proposed here are complementary rather than substitutable: regulation raises the security floor of the device itself, while this project addresses the in-life monitoring and response problem that persists regardless of how well-built the device was on the day it shipped.

These instruments raise the floor for newly manufactured devices, but they do not remediate the several billion units already deployed, nor do they address in-life detection and response. They therefore reinforce, rather than remove, the need for network-side defensive architectures capable of protecting devices that cannot protect themselves.

\subsection{Toward Autonomous, Self-Healing Defence}

Taken together, these conditions define the design space for this project. The endpoint population is large, growing, physically consequential and computationally incapable of self-defence. The threat is capable of sub-ten-minute, multi-terabit action assembled from precisely those endpoints, and the same population is independently and repeatedly targeted at the operator level, often by well-resourced adversaries. Detection latency is increasing rather than decreasing, and the human triage capacity that conventional architectures assume is already saturated by both alert volume and a persistent staffing shortfall. Regulation is tightening the baseline for future devices but leaves the installed base exposed, and does not itself provide an in-life response capability.

Securing such environments requires a shift from static, human-mediated, diagnostic defence toward dynamic architectures that detect, explain and remediate autonomously, at machine speed, while remaining constrained by the physical realities of the hardware they protect. That shift is the subject of this report.

\section{Problem Statement}

Traditional Intrusion Detection Systems (IDS) generally suffer from three critical limitations in modern IoT contexts. First, purely centralized cloud-based systems introduce unacceptable latency and bandwidth overhead, while isolated edge-based systems lack the processing power for complex threat classification. Second, most machine learning-based security tools operate as ``black boxes.'' When an anomaly is detected, network administrators are provided with an alert but lack the context to understand why the decision was made, leading to alert fatigue and delayed incident response. Finally, conventional frameworks are largely diagnostic rather than restorative; they can identify a breach but rely on manual human intervention to mitigate the threat and restore network integrity.

Each of these limitations is grounded in a distinct, measurable failure mode rather than a purely theoretical concern, and each is examined in turn below before the three are drawn together into the single operational gap that motivates this project.

\subsection{The Latency-Capability Trade-off in Centralized and Edge-Only Architectures}

The first limitation is architectural rather than algorithmic: wherever the intelligence of a defensive system is placed, that placement carries a cost that the threat model reviewed in Section 1.1 does not tolerate. A purely centralized design, in which raw or lightly filtered telemetry from every endpoint is streamed continuously to a cloud service for analysis, is attractive because it concentrates computational power where it is cheapest to provide; an unconstrained cloud instance can run a deep temporal model or a large ensemble classifier that no gateway could host. But this design pays for that capability in round-trip time and in bandwidth that a constrained network simply may not have. Every additional hop between the point of observation and the point of decision adds queuing delay, transmission delay, and processing delay, and in a fleet of many thousands of devices, the aggregate telemetry volume required to give a cloud classifier a complete view of every endpoint, at all times, competes directly with the operational traffic the network exists to carry. This is precisely the tension already identified in RFC 7228's characterisation of constrained devices and their gateways \cite{Bormann2014}: the very devices that most need protecting are the least able to absorb the overhead of continuous, high-fidelity reporting.

The opposite design, in which all detection logic is pushed to the edge and the cloud is dispensed with entirely, fails for the complementary reason. An edge gateway, even one considerably more capable than the Class 0 to Class 2 endpoints it protects, still operates under a real power, memory and compute budget, and cannot host the kind of high-capacity ensemble classifier needed to distinguish between attack types with the confidence that an operator would require before authorising an automated response. An edge-only detector can, at best, tell an operator that something is anomalous; it cannot reliably tell them what is anomalous, and the distinction matters operationally, because the correct remediation for a volumetric flood is not the same as the correct remediation for data exfiltration or for lateral reconnaissance. A system that cannot distinguish between these is left choosing between under-reacting to a serious intrusion and over-reacting to a benign traffic pattern, and either error carries a cost.

The severity of this trade-off is set by how little time is actually available. As established in Section 1.1.3, Cloudflare's own measurements found that 71\% of HTTP-layer and 89\% of network-layer DDoS attacks drawn from this same endpoint population concluded in under ten minutes \cite{Cloudflare2025Q3}, and the Aisuru-attributed record attacks it documented lasted 69 and 35 seconds respectively \cite{Cloudflare2025Q3, Cloudflare2026Q4}. Against this timescale, even a well-run security operations centre is working against a mismatched clock: industry benchmarking places an operationally acceptable mean time to respond, once an incident has already been detected and triaged by a human team, at roughly two to four hours for the full pipeline of acknowledgement, investigation and containment, and considerably longer for organisations without mature monitoring in place. A defensive pipeline whose fastest achievable human-mediated response measures in hours cannot be brought to bear on an attack that concludes in under a minute, regardless of how accurate its eventual classification proves to be; by the time telemetry has been forwarded, triaged, and acted upon through a conventional operational workflow, the incident under investigation has typically already ended. This is not a criticism of any particular tool or team; it is a structural mismatch between the tempo of the threat and the tempo of the response pipeline built to counter it, and it is the mismatch that motivates keeping an initial detection-and-response capability physically close to the endpoint rather than routing every decision through a distant, human-supervised tier.

\subsection{Opacity of Machine-Learning-Based Detection and the Erosion of Analyst Trust}

The second limitation concerns not whether a threat is detected, but whether the detection can be acted upon with confidence. Modern anomaly and intrusion detection increasingly depends on deep learning and large ensemble models precisely because these models achieve higher detection accuracy than earlier signature- or rule-based approaches; but that accuracy is bought at the cost of interpretability. A systematic literature review of explainable AI in cybersecurity, synthesising 169 peer-reviewed studies, reports that security specialists were able to fully understand the reasoning behind only 38\% of the alerts generated by contemporary neural-network-based detectors, and that roughly two-thirds of security experts are unwilling to act promptly on an automated recommendation when the system cannot explain why it was made \cite{Oriaro2025}. The same review found that moving from an opaque model to an inherently interpretable one, such as a decision tree, typically costs 15--20\% of detection accuracy \cite{Oriaro2025}-meaning that the choice between an accurate model and an explainable one is, at present, a genuine trade-off rather than a solved problem, and that simply discarding black-box models in favour of transparent ones is not a viable general solution.

This opacity does not merely inconvenience an analyst; it actively compounds the alert-fatigue and staffing problems already documented in Section 1.1.4. An alert that cannot explain itself forces the analyst who receives it to reconstruct, from raw features or from memory of similar past incidents, why the model considered the traffic suspicious, before any decision about remediation can even begin. In an environment where, per the ISC2 workforce data already cited, only around a third of organisations report having adequate security staffing \cite{ISC22025}, and where the 2025 SANS survey found false positives to be the leading operational complaint of detection teams \cite{SANS2025}, every additional minute an analyst spends decoding a model's reasoning is a minute not spent on the next alert in an already backlogged queue. Independent survey evidence supports the scale of the resulting distrust: broader human-factors research on AI-assisted security operations has found that a majority of analysts express open skepticism toward automated alerts and prefer a hybrid human-AI model over full automation specifically because they cannot verify an opaque system's reasoning under time pressure. The practical consequence is that an uninterpretable detector, however accurate in isolation, tends to be either ignored once alert volume exceeds what can be manually verified, or acted upon reflexively without verification, and neither outcome is acceptable in a system meant to authorise network-level remediation actions such as VLAN isolation or bandwidth throttling. An architecture that intends to act automatically on a model's output, rather than merely display that output to a human, must therefore treat explainability not as a usability nicety but as a precondition for safely delegating the decision at all-which is the requirement this project addresses by carrying SHAP-based feature attributions all the way through to the agent that ultimately authorises a remediation action, rather than stopping at a human-facing summary plot.

\subsection{The Manual-Intervention Bottleneck in Diagnostic Frameworks}

The third limitation is that, even where a threat is both detected and correctly explained, conventional frameworks still treat remediation as a separate, human-initiated step rather than as part of the same pipeline. This division has a measurable cost. Industry-reported benchmarks for mean time to respond-the interval from confirmed detection to full containment-place a commonly cited target range at two to four hours per incident for teams with a mature incident-response process, and considerably longer where the process is less mature or where escalation across multiple teams is required; the global averages implied by breach-lifecycle reporting are wider still, with Mandiant's most recent figures showing a median dwell time, from initial compromise to detection alone, of 14 days, before any remediation activity has begun \cite{Mandiant2026}. IBM's 2025 Cost of a Data Breach Report offers a direct illustration of what closing this gap is worth: organisations making extensive use of AI and automation in their security operations shortened their breach lifecycle by 80 days and reduced the average cost of an incident by USD 1.9 million relative to organisations using little or no automation \cite{IBM2025}. The saving is not incidental; it is a direct function of removing the human-mediated steps between detection and containment, which is exactly the step this project automates for the IoT context specifically.

The cost of manual mediation is amplified, rather than merely repeated, in the IoT setting because the underlying endpoint population is not compromised as an isolated, one-off event but is subject to sustained, repeated targeting. The utility-sector survey already discussed in Section 1.1.3 found that 80\% of operators who had been attacked at all had been attacked more than once within the same year, and that in the majority of cases the initial point of compromise was not the field device itself but an adjoining identity or network system, meaning that a single successful manual remediation does not close the underlying exposure \cite{Semperis2025}. A framework that requires a human operator to manually reconfigure a VLAN, adjust a firewall rule, or throttle bandwidth every time a device is compromised is therefore not solving the problem once; it is committing to repeating the same manual workflow indefinitely, against an adversary population that returns. Diagnostic-only frameworks are not wrong to stop at detection-identifying the attack vector correctly, as Chapter \ref{cha:methodology} discusses in the context of XGBoost-based classification, is a necessary precondition for any sound remediation decision-but stopping there leaves the most time-critical and most frequently repeated part of the incident-response workflow entirely dependent on a human being available, informed, and fast enough to act before the next occurrence of the same attack pattern.

\subsection{Synthesis: A Case for Autonomous, Closed-Loop Defence}

Taken together, these three limitations are not independent defects that could each be patched in isolation; they compound one another along the same timeline. An architecture that places its intelligence in the wrong tier loses the seconds-to-minutes window in which an IoT-scale attack actually occurs; a detector that cannot explain itself forces a staff-constrained human analyst to spend part of that already-scarce window on interpretation rather than action; and a framework that stops at diagnosis reintroduces, at the remediation step, exactly the delay that faster detection was meant to eliminate. Addressing any one of these limitations without the others would still leave the overall response pipeline too slow, too opaque, or too dependent on manual effort to match the tempo and persistence of the threat described in Section 1.1. This is the gap the proposed framework is designed to close: not by making any single stage marginally faster or more accurate, but by restructuring the pipeline itself so that detection, explanation and remediation occur as a single automated sequence, with a human operator positioned to supervise and escalate rather than to execute every step.

To address these challenges, this project proposes a novel, self-healing security framework for IoT networks that bridges lightweight edge computing with cognitive, cloud-based agentic Artificial Intelligence (AI). The proposed architecture introduces a multi-stage pipeline designed not only to detect and classify threats but to autonomously understand, explain, and remediate them.

\section{Objectives and Scope}

\subsection{Objectives}
\begin{itemize}
	\item Research the possibility of autonomic self-healing security of IoT systems by use of a layered intelligence architecture.
	\item Establish a security system that balances between fast response and dynamic thinking that will help to ensure that any threat can be contained in time and that the system will help ensure resiliency in the long term.
	\item Provide support to explainable and trustful decision-making during automated IoT security response by incorporating explainable AI procedures.
	\item Research the application of large language models to extrapolate to previously unknown patterns of attack, especially zero-day and evolving attacks.
	\item Create a self-directed learning strategy that facilitates the ongoing revision and enhancement of security policies within dynamic Internet of Things environments.
\end{itemize}

\subsection{Scope}

The research scope is formulated to provide practical relevance and feasibility of the experiment.

\begin{description}
	\item[System Scope] The framework focuses on networked IoT settings, such as smart home, industrial IoT (IIoT), and enterprise IoT implementations. Indirect protection of resource-constrained IoT endpoints is achieved by edge nodes and centralized layers of intelligence \cite{Bormann2014, Firouzi2025}.
	\item[Security Scope] The study targets network-based and behaviour-based attacks such as DDoS, botnet operation, scanning, protocol misuse, and normal traffic patterns \cite{Antonakakis2017, Jain2024, Cloudflare2025Q3}. Physical intrusions and mechanical interference are beyond the scope of this study.
	\item[AI Scope] The research takes into account classic ML-based known-attack detectors and LLM-based reasoning over higher-level logic and healing-strategy generation. The LLM is deployed in either an advisory or a controlled mode to reduce operational risk. Explainability mechanisms (e.g., local explanations and feature attribution) are incorporated to justify automated decisions and improve trustworthiness \cite{Lundberg2017, Oriaro2025, Tariq2025}.
	\item[Healing Scope] The actions that form part of healing are network isolation, rate limiting, credential reset, policy changes, and restricted reboot \cite{Jain2024}. Autonomous healing is limited by risk level, and high-impact actions are approved by a human-in-the-loop.
	\item[Deployment Scope] The framework is also intended to be modular and scalable, with and without cloud-assisted deployment. The use of federated or centralized learning is discussed as an extension, but is not the main priority of the implementation reported here \cite{McMahan2017, HernandezRamos2025}.
\end{description}


%%%%%%%%%%%%%%%%%%%%%%%%%%%%%%%%%%%%%%%%%%%%%%%%%%%%%%%%%%%%%%%%%%%%%%%%%%%%%%%%%%%%%%%%%%%%%%%%%%
%%%%%%%%%%% END OF FILE
%%%%%%%%%%%%%%%%%%%%%%%%%%%%%%%%%%%%%%%%%%%%%%%%%%%%%%%%%%%%%%%%%%%%%%%%%%%%%%%%%%%%%%%%%%%%%%%%%%